\section{Rational one-prime vanishing and actual residual matches}
\label{sec:rational-one-prime-residual-match}

\paragraph{Original coefficients and the remaining component.}
For an odd prime $\ell$, a finite extension $E/\mathbb Q_\ell$
and its original ring of integers $O$, put $D_O=E/O$.
For a continuous scalar action $a:G_{\mathbb Q}\to O^\times$,
write $\mathcal S^{\mathbb Q}_\ell(D_O(a))$ for the kernel of
continuous global $H^1$ restricted to the full local group at
$\ell$, to inertia at every other finite prime, and to the
real-place $H^1$. The real term is zero at odd $\ell$.
The full local condition at $\ell$ is essential.
All coefficient rings below are retained, including their
ramification and the specified embeddings of finite character values.

Let $c_\ell$ be the ordinary cyclotomic character.
Its reduction takes values in the prime field
$\mathbb F_\ell\subset O/\mathfrak m_O$.
Let $\tau_\ell:G_{\mathbb Q}\to\mathbb Z_\ell^\times$ be its
Teichm\"uller lift, with image $\mu_{\ell-1}$.
Thus $\tau_\ell$ is the finite Dirichlet character modulo $\ell$
characterized by
\[
 \overline{\tau_\ell(g)}=\overline{c_\ell(g)},\qquad
 \tau_\ell(g)^{\ell-1}=1.
 \tag{RQM.1}
\]
These values already lie in every such original $O$.
This definition fixes the exponent convention; no inverse
cyclotomic convention is substituted.

\begin{theorem}[AC-G195: original-integral rational residual matching]
\label{AC-G195}
Let $\ell$ be odd. If $a:G_{\mathbb Q}\to1+\mathfrak m_O$
is continuous and unramified at every finite prime other than $\ell$,
then
\[
                   \mathcal S^{\mathbb Q}_\ell(D_O(a))=0.
 \tag{RQM.2}
\]
In particular, let $\psi$ be a finite Dirichlet character with values
in this $O$, ramified only at $\ell$, and let $e\geq1$ satisfy
\[
                       \overline\psi=\overline{c_\ell}^{\,e}
                              \quad\hbox{on }G_{\mathbb Q}.
 \tag{RQM.3}
\]
Then $\mathcal S^{\mathbb Q}_\ell(D_O(\psi c_\ell^{-e}))=0$.
The finite parity is necessarily $\psi(-1)=(-1)^e$, so this is
exactly the noncritical rational component in (CP.3).

Keep the actual host, original $O_\ell$ and $b_\ell$ of (FC.1)--(FC.2),
with $K=\mathbb Q(\omega)$, $h\geq1$, and the distinct source primes
$q,d>7$ satisfying $q,d\equiv3\pmod4$. Let $T\geq1$ be an integer.
For $r\in\{q,d\}$, let $\lambda_r$ be the primitive quadratic
Dirichlet character modulo $r$, that is the Legendre character.
It is the character of the actual factor $\mathbb Q(\sqrt{-r})$.
The following internally disjoint families consist of actual
trivial-cubic triples $(\ell,\eta,e)$:
\[
 \begin{aligned}
 \mathfrak P_+(T)&=\{\ell\leq T:\ell\geq7\text{ prime},\
                  \ell\equiv1\pmod3,\ \ell\ne7\},\\
 Y_+(T)&=\{(\ell,1,e):\ell\in\mathfrak P_+(T),\
                        1\leq e<T,\ (\ell-1)\mid e\},\\
 Y_7(T)&=\{(7,\tau_7^e|_{G_K},e):1\leq e<T\},\\
 Y_r(T)&=\{(r,\lambda_r|_{G_K},e):1\leq e<T,\
                    e\equiv(r-1)/2\pmod{r-1}\}\quad(r=q,d),\\
 Y(T)&=Y_+(T)\sqcup Y_7(T)\sqcup Y_q(T)\sqcup Y_d(T).
 \end{aligned}                                                    \tag{RQM.4}
\]
The finite factors not displayed in each label are trivial.
For every triple in $Y(T)$, the rational extension displayed in
(RQM.4) is its unique noncritical extension, and its rational
full strict group is zero. On the same original coefficient there
is an exact identification
\[
 \mathcal S_{\ell,\eta}(e)
       \simeq\mathcal S^{\mathbb Q}_{\ell,
                                  \psi_{\eta,e}^{\rm cr}}(e).
 \tag{RQM.5}
\]
This does not assert that the group in (RQM.5) is zero.
For any subset $Y'\subseteq Y(T)$ and $T\geq2$, the already proved
critical aggregate estimate gives
\[
  \sum_{(\ell,\eta,e)\in Y'}
                \frac{\log\#\mathcal S_{\ell,\eta}(e)}{b_\ell}
          \leq\mathcal B_{\rm cr}(F,T)\leq C\mathcal P(F,T),
 \tag{RQM.6}
\]
with exactly the $F$, $\mathcal P$ and absolute constant scope of
AC-G191. The original nonnegative remainder (CP.3) satisfies
$\mathcal R_J=\mathcal R_{J\setminus Y(T)}$ for every retained set $J$.

The exact count, with $N_T=T-1$, is
\[
 \begin{split}
 J_+(T)&=\sum_{\ell\in\mathfrak P_+(T)}
                      \left\lfloor\frac{T-1}{\ell-1}\right\rfloor,\\
 A_r(T)&=\left\lfloor
       \frac{T-1+(r-1)/2}{r-1}\right\rfloor\quad(r=q,d),\\
 |Y(T)|&=N_T+J_+(T)+A_q(T)+A_d(T).
 \end{split}                                                    \tag{RQM.7}
\]
In particular the auxiliary counts can be zero; their presence
below the actual cutoff is not assumed. The coefficient-seven
list includes the trivial label when $6\mid e$, so that prime is
excluded from $Y_+$ to avoid counting those triples twice.
No individual positive derived index is asserted zero.
\end{theorem}

\begin{proof}
\proofstep{1. The rational local decomposition group generates the pro-$\ell$ quotient.}
Let $\Gamma$ be the Galois group of the maximal extension of
$\mathbb Q$ unramified at all finite primes other than $\ell$,
and let $P=\Gamma^{(\ell)}$ be its maximal pro-$\ell$ quotient.
The image of $G_{\mathbb Q_\ell}$ in $P$ is compact and hence closed.
If it were proper, some finite $\ell$-group quotient $Q$ of $P$
would have proper local image. Every proper subgroup of a finite
$\ell$-group is contained in a normal subgroup of index $\ell$.
For completeness, a maximal proper subgroup is normal of that
index by induction using the nontrivial center: if the center is
not contained in it, their product is the whole group; otherwise
pass to the quotient by the center.
Thus a proper local image would give a nontrivial cyclic extension
of $\mathbb Q$ of degree $\ell$, unramified away $\ell$ and split
at $\ell$. It is unramified at every finite prime.
It is also split at the real place, because a real decomposition
group has order at most two and $\ell$ is odd.
Ordinary Hilbert class-field reciprocity would make its Galois
group a quotient of $\operatorname{Cl}(\mathbb Q)$.
Every ideal of $\mathbb Z$ is principal, so this class group is
trivial. This is a contradiction. Hence
\[
                      G_{\mathbb Q_\ell}\longrightarrow P
                                 \quad\hbox{is surjective}.
 \tag{RQM.8}
\]
This is the real-field odd-prime form of the elementary argument
proved in AC-G193, specialized to $\mathbb Q$.

\proofstep{2. The finite image includes the original local witness.}
Let $z:G_{\mathbb Q}\to D_O(a)$ represent a strict class.
Its image is finite by compactness and discreteness.
Full local zero gives an element $x\in D_O$ such that
$z(g)=(a(g)-1)x$ for all $g\in G_{\mathbb Q_\ell}$.
Choose a uniformizer $\pi$ of the original $O$ and $m$ large
enough that both the entire image of $z$ and this witness $x$
lie in $D_O[\pi^m]$.
The image $U_m$ of $1+\mathfrak m_O$ on this finite module is
an $\ell$-group, as is its additive group. The cocycle identity
makes
\[
 g\longmapsto\bigl(z(g),a(g)|_{D_O[\pi^m]}\bigr)
       \quad\hbox{a homomorphism into }D_O[\pi^m]\rtimes U_m.
 \tag{RQM.9}
\]
At every finite prime other than $\ell$, the coefficient is
unramified. A zero inertia cohomology class therefore means
$z(i)=0$ on that inertia group, since every inertia coboundary
is zero. The finite homomorphism (RQM.9) kills their closed
normal closure. It factors through $\Gamma$, and then through
$P$ because its image is an $\ell$-group.
By (RQM.8), for every global element there is a local element
with the same image in (RQM.9). Their actions on the chosen
witness agree, because the witness belongs to that finite module.
Its local coboundary identity therefore holds globally.
This proves (RQM.2) on the unchanged $O$.

For the coefficient in (RQM.3), the full scalar
$\psi c_\ell^{-e}$ lies in $1+\mathfrak m_O$ and is unramified
outside $\ell$, so the result applies. Evaluation at complex
conjugation gives $\psi(-1)(-1)^e\equiv1\pmod{\mathfrak m_O}$.
Both sides of the finite sign lie in $\{1,-1\}$, whose reductions
are distinct at odd $\ell$. Thus $\psi(-1)=(-1)^e$ exactly.
The real-place $H^1$ is zero because two is invertible on $D_O$.

\proofstep{3. The exact eligible finite characters.}
For the trivial rational character, (RQM.3) holds exactly when
$(\ell-1)\mid e$, since the residual cyclotomic character of
$G_{\mathbb Q}$ is surjective onto $\mathbb F_\ell^\times$.
For every split prime in $Y_+$ these are the stated weights.
No assertion about the other rational component is used here.

At $\ell=7$, the six finite characters modulo seven all have
values of order dividing six. Reduction at seven is injective
on these roots of unity. Consequently exactly one such character
has reduction $\overline{c_7}^{\,e}$: it is $\tau_7^e$.
It depends on $e$ modulo six, and is trivial precisely when
$6\mid e$. The actual host contains $\mathbb Q(\zeta_7)$ and
its original coefficient order has exponent $3\cdot2^h$, divisible
by six. Thus all six choices already occur over the specified
$O_7$, with its actual embedding. No new finite character is
adjoined to the host.

Here is a direct identification for the auxiliary factors.
For any odd prime $r$, put $\lambda_r(a)=(a/r)$, with
$\lambda_r(0)=0$, and let $z$ be a primitive $r$th root of unity.
The Gauss sum $g_r=\sum_{a\bmod r}\lambda_r(a)z^a$ satisfies
\[
 \sigma_b(g_r)=\lambda_r(b)g_r,\qquad
 \overline{g_r}=\lambda_r(-1)g_r,\qquad
                         g_r\overline{g_r}=r.
 \tag{RQM.10}
\]
The first two identities follow by substitution. For the third,
write $a=tb$ in the double sum for $g_r\overline{g_r}$, with
$b\ne0$. The inner sum over $b$ is $r-1$ at $t=1$ and $-1$
otherwise; using $\sum_t\lambda_r(t)=0$ gives $r$.
It follows that $g_r^2=\lambda_r(-1)r$ and that the fixed field
of $\ker\lambda_r$ is $\mathbb Q(\sqrt{\lambda_r(-1)r})$.
For $r=q,d\equiv3\pmod4$ this is exactly the actual
$\mathbb Q(\sqrt{-r})$ factor, of primitive conductor $r$.
Euler's criterion gives
\[
                 \lambda_r=\tau_r^{(r-1)/2},
   \qquad\overline{\lambda_r}=\overline{c_r}^{\,(r-1)/2}.
 \tag{RQM.11}
\]
Both values in the first identity are the same sign $\pm1$;
no coefficient extension is required.
Thus (RQM.3) holds at precisely the auxiliary weights in
(RQM.4). The notation $\lambda_r$ denotes this conductor-$r$
Legendre character, not the different character of
$\mathbb Q(\sqrt r)$ of conductor $4r$.

All the displayed characters factor through the abelian field
$F_0$ of (CP.1), with trivial Kummer cubic factor. The coprime
conductor factors in (FG.11) give $F_0\cap K=\mathbb Q$.
Their restrictions to $G_K$ are therefore distinct whenever
the rational characters are distinct. The auxiliary factors
have values $\pm1$, and the coefficient-seven factors have
values in $\mu_6$, already in the original $O_\ell$ of (FC.1).
The finite character at every unspecified factor is trivial.
This proves that (RQM.4) uses actual labels and actual completions.

\proofstep{4. Full strict groups and the two rational extensions.}
Let $\psi$ be any displayed rational character and
$\eta=\psi|_{G_K}$. Its two extensions are $\psi$ and
$\psi\chi_{-3}$. The parity just proved makes $\psi$ noncritical
and $\psi\chi_{-3}$ critical in the convention of (CP.1)--(CP.3).
On the original $O_\ell$ the rank-two induced lattice splits:
\[
 \operatorname{Ind}_{G_K}^{G_{\mathbb Q}}O_\ell(\eta)
        \simeq O_\ell(\psi)\oplus O_\ell(\psi\chi_{-3}).
 \tag{RQM.12}
\]
The two-coset evaluation matrix has determinant a unit times two,
which is invertible for every coefficient prime used here.
Tensor with the original $E_\ell/O_\ell$ and apply the same Tate
twist on both sides. Shapiro gives the corresponding global
cohomology isomorphism.

This identifies the full strict kernels as well.
Local Shapiro at $\ell$ includes every place of $K$ above $\ell$,
including both places at split primes. At any other prime,
Mackey applied to inertia gives the actual inertia subgroups in
$K$. At the ramified prime three its intersection with $G_K$
is the inertia group of $K$ and there is one inertia double coset.
At an unramified prime the two inertia components are either
the two split places or conjugate components at the single inert
place; their simultaneous vanishing is exactly the original
inertia condition. At infinity both rational $H^1$ groups vanish
at odd $\ell$, consistently with induction from a complex place.
This proves the exact direct sum of full strict groups.

The global group in (FC.2) may be used without changing this
identification. Its localization set contains every ramification
prime of these coefficients. Inflation embeds its cohomology
into full global cohomology; outside the localization set an
inflated cocycle already kills inertia. Conversely, a full strict
cocycle kills all inertia outside that set, where the coefficient
is unramified, and its finite semidirect image factors through
the same localization quotient. All additional host primes
remain subject to the original inertia-zero condition.
We therefore obtain (RQM.5) from (RQM.2), on precisely the
original $O_\ell$ and with the normalization $b_\ell$ unchanged.

\proofstep{5. Uniform aggregate and exact count.}
Each critical summand in (RQM.5) is one of the nonnegative terms
of the full all-prime sum (CP.2). Different triples in any $Y'$
give different such terms: restriction of a rational character
recovers its original $\eta$ and its coefficient and weight are
retained. Thus (RQM.6) follows from AC-G191 without a multiplier
for the number of primes, auxiliary factors or labels.
Every removed term of (CP.3) is a zero noncritical term, proving
$\mathcal R_J=\mathcal R_{J\setminus Y(T)}$.

For a fixed $\ell\in\mathfrak P_+(T)$, the number of positive
multiples of $\ell-1$ below $T$ is the summand in $J_+(T)$.
For $\ell>T$ it is zero, which explains the finite prime range.
At coefficient seven there is one, and only one, eligible label
at each weight, giving $N_T$ terms. There are no coefficient-seven
terms in $Y_+$, so the trivial-label overlap has been removed.
For $r=q,d$ write $a=(r-1)/2$. The auxiliary weights are
$a,a+2a,a+4a,\ldots$; counting those at most $T-1$ gives
$\lfloor(T-1+a)/(2a)\rfloor=A_r(T)$, including zero when $T\leq a$.
The two auxiliary coefficient primes are distinct and exceed seven.
Their finite labels are nontrivial, so they cannot coincide with
any trivial-label term in $Y_+$, even when that coefficient is split.
These observations prove the internal disjointness and (RQM.7).
\end{proof}

\paragraph{Actual-window and source boundary.}
The estimate (RQM.6) is pointwise for every actual pair, page,
coefficient and weight range. At $T=\lfloor u^{399}\rfloor$ its
cost has the already proved polynomial envelope (CP.5).
A separate composition with the current omitted set must retain
one full (FL.4) local-error majorant and the identical common $W$.
This theorem does not split a positive derived index, declare
that an original-$K$ strict group is zero, or assert strict numerical
decrease of a remainder. The mixed cubic terms and all other
noncritical rational terms are retained.

The initial rational one-prime route was proposed by the parent
reviewer. The proof above repeats the needed finite-group and
finite-cocycle argument completely. Its deep input is ordinary
Hilbert class-field reciprocity, in Sharifi, \emph{Iwasawa Theory},
Theorem1.1.11 and its preceding definition; the complete applied
maximal-$S$ extension lemma is Lemma1.4.2 and Definition1.4.3.
The real odd-prime version and full-local witness are also proved
in AC-G193. The actual host and coefficient dictionaries are the
complete applied (FG.11), (FC.1)--(FC.2) and (CP.1), and the sole
quantitative input is the independently checked AC-G191.
No new class computation, narrow-class assumption, regulator
estimate or conjectural zero statement is used.
